% Generated by scripts/make_paper_assets.py. Do not edit by hand.
\begin{table}[t]
\centering
\small
\begin{tabular}{lcc}
\toprule
Split protocol & cesnet\_quic22 & fiveg\_traffic \\
\midrule
\multicolumn{3}{l}{\emph{FlowCon-X (prototype head)}} \\
Random flow & 0.781 $\pm$ 0.006 & 0.991 $\pm$ 0.000 \\
Session-disjoint & \textsc{todo} & 0.522 $\pm$ 0.139 \\
Temporal & 0.781 $\pm$ 0.005 & \textsc{todo} \\
Server-disjoint & 0.588 $\pm$ 0.007 & \textsc{todo} \\
\midrule
\multicolumn{3}{l}{\emph{AppScanner (2016), identical splits}} \\
Random flow & 0.768 & 0.996 \\
Session-disjoint & 0.772 & 0.640 \\
Temporal & 0.761 & 0.255 \\
Server-disjoint & 0.566 & 0.711 \\
\bottomrule
\end{tabular}
\caption{Macro-F1 under four split protocols, every method on the identical committed splits. The random-flow row is the protocol standard in the literature and is reported only as a contrast: correlated flows from one capture session appear on both sides of it. \textbf{The size of that contrast is dataset-dependent}, which is itself the result: on a corpus of single-session captures a flow's capture is nearly its label, while on a backbone corpus whose flows are already independent the protocol makes little difference. The caption deliberately does not assert a direction; read the columns.}
\label{tab:split-contrast}
\end{table}
